% Partie : applied-modal-logics
% Chapitre : epistemic-logic
% Section : public-announcement-logic-lang

\documentclass[../../../include/open-logic-section]{subfiles}

\begin{document}

\olfileid[fr]{aml}{el}{pal}

\olsection{Logique des annonces publiques}

Les logiques épistémiques dynamiques permettent de représenter les
changements de connaissance des agents au fil du temps ou lorsqu'ils
reçoivent de nouvelles informations. Plusieurs d'entre elles représentent
ces changements par des \emph{événements} informationnels ou des
\emph{mises à jour}. La mise à jour la plus élémentaire est une annonce
publique : une formule est annoncée de manière véridique et tous les
agents assistent ensemble à cette annonce. Pour la représenter, nous
étendons le langage comme suit.

\begin{defn}
Soit $G$ un ensemble de symboles d'agents. Le langage de base de la
logique épistémique à plusieurs agents avec annonces publiques contient :\footnote{L'inventaire de la source omet l'équivalence matérielle, pourtant admise par la clause de formation qui suit ; elle est rétablie ici.}
\begin{enumerate}
  \tagitem{prvFalse}{La constante propositionnelle de !!{falsity}~$\lfalse$.}{}
  \tagitem{prvTrue}{La constante propositionnelle de !!{truth}~$\ltrue$.}{}
  \item Un ensemble !!{denumerable} de !!{propositional variable}s : $\Obj
    p_0$, $\Obj p_1$, $\Obj p_2$, \dots
  \item Les connecteurs propositionnels : \startycommalist
  \iftag{prvNot}{\ycomma $\lnot$ (négation)}{}%
  \iftag{prvAnd}{\ycomma $\land$ (conjonction)}{}%
  \iftag{prvOr}{\ycomma $\lor$ (disjonction)}{}%
  \iftag{prvIf}{\ycomma $\lif$ (!!{conditional})}{}%
  \iftag{prvIff}{\ycomma $\liff$ (!!{biconditional})}{}.
  \item L'opérateur de connaissance $\Knows_a$, pour chaque $a \in G$.
  \item L'opérateur d'annonce publique $[!B]$, où $!B$ est !!a{formula}.
\end{enumerate}
\end{defn}

L'opérateur d'annonce publique se comporte comme un opérateur de
nécessité ; la définition inductive du langage en tient compte :

\begin{defn}
  Les \emph{!!{formula}s} du langage épistémique avec
  annonces publiques sont définies inductivement comme suit :
  \begin{enumerate}
  \tagitem{prvFalse}{$\lfalse$ est une !!{formula} atomique.}{}

  \tagitem{prvTrue}{$\ltrue$ est une !!{formula} atomique.}{}

  \item Toute variable propositionnelle $\Obj p_i$ est une
  !!{formula} (atomique).

  \tagitem{prvNot}{Si $!A$ est !!a{formula}, alors $\lnot !A$ est
  !!a{formula}.}{}

  \tagitem{prvAnd}{Si $!A$ et $!B$ sont des !!{formula}s, alors $(!A \land
  !B)$ est !!a{formula}.}{}

  \tagitem{prvOr}{Si $!A$ et $!B$ sont des !!{formula}s, alors $(!A \lor !B)$
  est !!a{formula}.}{}

  \tagitem{prvIf}{Si $!A$ et $!B$ sont des !!{formula}s, alors $(!A \lif !B)$
  est !!a{formula}.}{}

  \tagitem{prvIff}{Si $!A$ et $!B$ sont des !!{formula}s, alors $(!A \liff !B)$
  est !!a{formula}.}{}

  \item Si $!A$ est !!a{formula} et $a \in G$, alors $\Knows_a !A$ est
  !!a{formula}.

  \item Si $!A$ et $!B$ sont des !!{formula}s, alors $[!A] !B$ est !!a{formula}.

  \tagitem{limitClause}{Rien d'autre n'est !!a{formula}.}{}
\end{enumerate}
\end{defn}

La lecture voulue de la !!{formula} $[!A] !B$ est : « Après une annonce
véridique de $!A$, $!B$~est vraie. » Il est parfois utile de parler
également de connaissance commune dans le contexte des annonces
publiques. On peut donc étendre aussi le langage par l'opérateur de
connaissance commune~$\CKnows_G
!A$.

\end{document}
